% Placeholder cover letter. Compile with XeLaTeX or LuaLaTeX.
% Matching Markdown: templates/application/cover-letter.md.
\documentclass[a4paper,10pt]{article}

\usepackage[british]{babel}
\usepackage{calc}
\usepackage{fontspec}
\usepackage{geometry}
\usepackage{microtype}
\usepackage{xcolor}

\newgeometry{left=15mm,right=15mm,top=15mm,bottom=15mm}
\newlength{\cvcolumngapwidth}
\setlength{\cvcolumngapwidth}{3.5mm}
\newlength{\cvleftcolumnwidth}
\setlength{\cvleftcolumnwidth}{44.5mm}
\newlength{\cvrightcolumnwidth}
\setlength{\cvrightcolumnwidth}{\textwidth-\cvleftcolumnwidth-\cvcolumngapwidth}
\setlength{\parindent}{0mm}
\setlength{\parskip}{2.5mm plus 0.5mm minus 0.5mm}

\defaultfontfeatures{Ligatures=TeX}
\newif\ifcvroboto
\cvrobototrue
\IfFontExistsTF{Roboto Light}{}{\cvrobotofalse}
\IfFontExistsTF{Roboto Medium}{}{\cvrobotofalse}
\ifcvroboto
  \newfontfamily{\cvnamefont}{Roboto Medium}
  \newfontfamily{\cvsectionfont}{Roboto Medium}
  \setmainfont{Roboto Light}
\else
  \newfontfamily{\cvnamefont}{lmsans10-regular.otf}
  \newfontfamily{\cvsectionfont}{lmsans10-regular.otf}
  \setmainfont[BoldFont=lmroman10-bold.otf,ItalicFont=lmroman10-italic.otf,BoldItalicFont=lmroman10-bolditalic.otf]{lmroman10-regular.otf}
\fi

\definecolor{cvnamecolor}{HTML}{000000}
\definecolor{cvsectioncolor}{HTML}{000000}
\definecolor{cvmaincolor}{HTML}{000000}
\definecolor{cvrulecolor}{HTML}{000000}
\color{cvmaincolor}

\newcommand{\cvnamestyle}[1]{{\Large\cvnamefont\textcolor{cvnamecolor}{#1}}}
\newcommand{\cvsectionstyle}[1]{{\normalsize\cvsectionfont\textcolor{cvsectioncolor}{#1}}}

\newlength{\cvafterheaderskipamount}
\setlength{\cvafterheaderskipamount}{7mm plus 1mm minus 1mm}
\newlength{\cvaftersectionskipamount}
\setlength{\cvaftersectionskipamount}{2mm plus 0.5mm minus 0.5mm}

\newcommand{\cvheader}[2]{
    \begin{minipage}[t]{\cvleftcolumnwidth}
        \raggedleft\cvnamestyle{#1}
    \end{minipage}%
    \hspace{\cvcolumngapwidth}%
    \begin{minipage}[t]{\cvrightcolumnwidth}
        #2
    \end{minipage}
    \par\vspace{\cvafterheaderskipamount}
}
\newcommand{\cvsection}[1]{
    \begin{minipage}[t]{\cvleftcolumnwidth}
        \raggedleft\cvsectionstyle{#1}
    \end{minipage}%
    \hspace{\cvcolumngapwidth}%
    \begin{minipage}[t]{\cvrightcolumnwidth}
        \textcolor{cvrulecolor}{\rule{\cvrightcolumnwidth}{0.3mm}}
    \end{minipage}
    \par\vspace{\cvaftersectionskipamount}
}
\newcommand{\cvlettertext}[1]{
    \begin{minipage}[t]{\cvleftcolumnwidth}
        \mbox{}
    \end{minipage}%
    \hspace{\cvcolumngapwidth}%
    \begin{minipage}[t]{\cvrightcolumnwidth}
        \setlength{\parindent}{0mm}
        \setlength{\parskip}{2.5mm plus 0.5mm minus 0.5mm}
        #1
    \end{minipage}
    \par\vspace{4mm}
}

\pagestyle{empty}

\begin{document}

\cvheader{[Full name]}{
    [City, country]\par
    [Optional phone]\par
    [Email]\par
    [Professional link]
}

\cvsection{WHY THIS COMPANY?}

\cvlettertext{
[Company] interests me because [specific, verified aspect of its work, product, or mission]. I am applying for [Role] because [connection between the role and my confirmed career interests].

[A second reason based on the role posting or company research]. I would welcome the opportunity to [contribution relevant to the role, without implying experience I have not documented].
}

\cvsection{WHY ME?}

\cvlettertext{
I bring [relevant strengths supported by my career records]. These connect to your need for [requirement from the role posting].

In [Role or project] at [Organisation], I [specific personal contribution]. [Supported outcome, keeping its scope and any measurement accurate]. This experience is relevant to [requirement or problem in the target role].

I also [second supported contribution or relevant way of working]. [Concrete evidence from a documented experience].

I would welcome the opportunity to discuss how my experience could contribute to [Company] as [Role].
}

\end{document}
